\documentclass[12pt]{article}
\usepackage[utf8]{inputenc}
\usepackage[T1]{fontenc}
\usepackage{hyperref}
\title{Tomadino Language Resources for K-12 Curriculum Development}
\author{Compiled by Claude Sonnet 5 (claude-sonnet-5) for Andr\'as Kornai}
\date{}
\begin{document}
\maketitle
\section{Introduction}
Tomadino (ISO 639-3 tdi; Glottocode toma1248) is an endangered Austronesian language of the Bungku--Tolaki group in Central Sulawesi, Indonesia. It is reported at about 600 speakers. Speakers are shifting to the larger Bungku language. Joshua Project gives a much lower population figure of about 100, and the sources conflict.
\section{Primary Educational Resources}
None found. The language has no literacy tradition.
\section{Linguistic and Reference Materials}
\begin{itemize}
\item Wikipedia overview: \url{https://en.wikipedia.org/wiki/Tomadino_language}
\item Austronesian Basic Vocabulary Database wordlist: \url{https://abvd.eva.mpg.de/austronesian/language.php?id=914}
\item Verbix language page: \url{https://docs.verbix.com/Languages/Tomadino}
\end{itemize}
The search summary named Mead's 1999 Pacific Linguistics volume on the Bungku--Tolaki languages as the main source. No URL for it appeared in the results.
\section{Religious and Community Resources}
Joshua Project people-group page: \url{https://joshuaproject.net/people_groups/15546}
\section{Digital and Audio Resources}
Endangered Languages Project entry: \url{https://endangeredlanguages.com/lang/2796}
\section{Curriculum Development Context}
Very little exists. Any curriculum work would first need documentation (wordlists, texts, recordings) and would have to consider the nearby Bungku language, which speakers are shifting to.
\end{document}